\documentclass{article}

% Language setting
% Replace `english' with e.g. `spanish' to change the document language
\usepackage[english]{babel}

% Set page size and margins
% Replace `letterpaper' with `a4paper' for UK/EU standard size
\usepackage[letterpaper,top=2cm,bottom=2cm,left=3cm,right=3cm,marginparwidth=1.75cm]{geometry}

% Useful packages
\usepackage{amsmath}
\usepackage{graphicx}
\usepackage[colorlinks=true, allcolors=blue]{hyperref}

\title{?}
\author{}

\begin{document}
	
	\section{Dataset Description}
	
	The dataset used in this study consists of multi-channel EEG recordings from mice treated with a chemoconvulsant to induce spontaneous seizures. The data were originally published in Lisgaras and Scharfman (2022)\footnote{\url{https://doi.org/10.1016/j.nbd.2022.105637}} and are shared for research purposes through a Hackathon organized by the ABQ Chapter of the ASA.
	
	\subsection{Data Structure and Format}
	
	Each recording session is stored as a binary \texttt{.dat} file containing 16-bit integers sampled at 2000 Hz across four channels. These channels correspond to:
	
	\begin{itemize}
		\item CH1: Ipsilateral Hippocampus
		\item CH2: Ipsilateral Neocortex
		\item CH3: Contralateral Hippocampus
		\item CH4: Contralateral Neocortex
	\end{itemize}
	
	Each \texttt{.dat} file is accompanied by a \texttt{.txt} annotation file that records event times, including seizure onsets and offsets.
	
	\subsection{Dataset File Summary}
	
	The following files were found in the dataset directory, along with their corresponding shapes (rows $\times$ channels):
	
	\begin{itemize}
		\item \texttt{AC75a-5\_DOB\_072519\_TS\_2020-03-23\_17\_30\_04\_allCh.dat}: shape = (172700000, 4)
		\item \texttt{AC75a-5\_DOB\_072519\_TS\_2020-03-24\_17\_30\_04\_allCh.dat}: shape = (172800000, 4)
		\item \texttt{AC75a-5\_DOB\_072519\_TS\_2020-03-25\_17\_30\_04\_allCh.dat}: shape = (172800000, 4)
		\item \texttt{AC75a-5\_DOB\_072519\_TS\_2020-03-27\_17\_30\_04\_allCh.dat}: shape = (172800000, 4)
		\item \texttt{AC75a-5\_DOB\_072519\_TS\_2020-03-28\_17\_30\_04\_allCh.dat}: shape = (172800000, 4)
		\item \texttt{AC75a-5\_DOB\_072519\_TS\_2020-03-29\_17\_30\_04\_allCh.dat}: shape = (172800000, 4)
		\item \texttt{AC75a-5\_DOB\_072519\_TS\_2020-03-30\_17\_30\_04\_allCh.dat}: shape = (172800000, 4)
		\item \texttt{AC75a-5\_DOB\_072519\_TS\_2020-03-31\_17\_30\_04\_allCh.dat}: shape = (172800000, 4)
		\item \texttt{AC75a-5\_DOB\_072519\_TS\_2020-04-02\_17\_30\_04\_allCh.dat}: shape = (172800000, 4)
		\item \texttt{AC75a-5\_DOB\_072519\_TS\_2020-04-04\_17\_30\_04\_allCh.dat}: shape = (172800000, 4)
		\item \texttt{AC75a-5\_DOB\_072519\_TS\_2020-04-06\_17\_30\_04\_allCh.dat}: shape = (172800000, 4)
	\end{itemize}
	
	
	The filenames encode metadata about the subject and the recording time:
	\begin{itemize}
		\item \texttt{AC75a-5}: Mouse ID
		\item \texttt{DOB\_072519}: Date of birth (July 25, 2019)
		\item \texttt{TS\_YYYY-MM-DD\_HH\_MM\_SS}: Start time of recording session
	\end{itemize}
	
	\subsection{Excluded File}
	
	The first recording file (\texttt{2020-03-23}) was excluded from further analysis. While the file contains 172{,}700{,}000 samples, a full 24-hour recording at 2000 Hz should contain exactly 172{,}800{,}000 samples. This discrepancy suggests that part of the data is missing. Since the binary \texttt{.dat} files do not contain embedded timestamps, it is unclear which segment is missing. As a result, the signal cannot be reliably aligned with the annotations in the corresponding \texttt{.txt} file, and the file was not used in our analysis.
	
	\subsection{Annotations}
	
	Each \texttt{.dat} file is associated with a corresponding \texttt{.txt} annotation file. These annotation files contain time-stamped events observed during the recording session. The entries include various behavioral or experimental events, with particular focus on seizures.
	
	The annotation files are formatted as tables, with each row representing an event. Key columns include:
	\begin{itemize}
		\item \textbf{Start Time} and \textbf{End Time}: The clock time at which the event occurred, with millisecond resolution.
		\item \textbf{Time From Start}: The time in seconds since the beginning of the recording.
		\item \textbf{Channel}: Typically marked as \texttt{ALL}, indicating that the event is global and not channel-specific.
		\item \textbf{Annotation}: A label describing the event, such as \texttt{Seizure starts}, \texttt{Seizure ends}, \texttt{Blinking}, or \texttt{Lights OFF}.
	\end{itemize}
	
	
	To align annotations with the EEG signals, we rely on the assumption that the first data sample in the corresponding \texttt{.dat} file maps to the \texttt{Start Time} of the recording session. This allows us to create a time index for the EEG signal and mark seizure intervals precisely for visualization and prediction tasks.
	
	
	\section{Seizure Highlight Visualizations}
	
	The following figures show the EEG recordings for each file (24-hour session). Each plot displays the four-channel EEG signal with seizure periods highlighted in red.
	
	\begin{figure}[h!]
		\centering
		\includegraphics[width=\textwidth]{C:/Users/Sina.Mokhtar.XLSCIENTIFIC/Documents/Problems/Copula_AD/figures/AC75a-5_DOB_072519_TS_2020-03-24_17_30_04_allCh.png}
		\caption{EEG plot from 2020-03-24 17:30:04 to 2020-03-25 17:30:04.}
	\end{figure}
	
	% \begin{figure}[h!]
		%     \centering
		%     \includegraphics[width=\textwidth]{figures/AC75a-5_DOB_072519_TS_2020-03-25_17_30_04_allCh.png}
		%     \caption{EEG plot from 2020-03-25 17:30:04 to 2020-03-26 17:30:04.}
		% \end{figure}
	
	% \begin{figure}[h!]
		%     \centering
		%     \includegraphics[width=\textwidth]{figures/AC75a-5_DOB_072519_TS_2020-03-27_17_30_04_allCh.png}
		%     \caption{EEG plot from 2020-03-27 17:30:04 to 2020-03-28 17:30:04.}
		% \end{figure}
	
	\begin{figure}[h!]
		\centering
		\includegraphics[width=\textwidth]{C:/Users/Sina.Mokhtar.XLSCIENTIFIC/Documents/Problems/Copula_AD/figures/AC75a-5_DOB_072519_TS_2020-03-28_17_30_04_allCh.png}
		\caption{EEG plot from 2020-03-28 17:30:04 to 2020-03-29 17:30:04.}
	\end{figure}
	
	% \begin{figure}[h!]
		%     \centering
		%     \includegraphics[width=\textwidth]{C:/Users/Sina.Mokhtar.XLSCIENTIFIC/Documents/Problems/Copula_AD/figures/AC75a-5_DOB_072519_TS_2020-03-29_17_30_04_allCh.png}
		%     \caption{EEG plot from 2020-03-29 17:30:04 to 2020-03-30 17:30:04.}
		% \end{figure}
	
	% \begin{figure}[h!]
		%     \centering
		%     \includegraphics[width=\textwidth]{C:/Users/Sina.Mokhtar.XLSCIENTIFIC/Documents/Problems/Copula_AD/figures/AC75a-5_DOB_072519_TS_2020-03-30_17_30_04_allCh.png}
		%     \caption{EEG plot from 2020-03-30 17:30:04 to 2020-03-31 17:30:04.}
		% \end{figure}
	
	\begin{figure}[h!]
		\centering
		\includegraphics[width=\textwidth]{C:/Users/Sina.Mokhtar.XLSCIENTIFIC/Documents/Problems/Copula_AD/figures/AC75a-5_DOB_072519_TS_2020-03-31_17_30_04_allCh.png}
		\caption{EEG plot from 2020-03-31 17:30:04 to 2020-04-01 17:30:04.}
	\end{figure}
	
	% \begin{figure}[h!]
		%     \centering
		%     \includegraphics[width=\textwidth]{C:/Users/Sina.Mokhtar.XLSCIENTIFIC/Documents/Problems/Copula_AD/figures/AC75a-5_DOB_072519_TS_2020-04-02_17_30_04_allCh.png}
		%     \caption{EEG plot from 2020-04-02 17:30:04 to 2020-04-03 17:30:04.}
		% \end{figure}
	
	% \begin{figure}[h!]
		%     \centering
		%     \includegraphics[width=\textwidth]{C:/Users/Sina.Mokhtar.XLSCIENTIFIC/Documents/Problems/Copula_AD/figures/AC75a-5_DOB_072519_TS_2020-04-04_17_30_04_allCh.png}
		%     \caption{EEG plot from 2020-04-04 17:30:04 to 2020-04-05 17:30:04.}
		% \end{figure}
	
	% \begin{figure}[h!]
		%     \centering
		%     \includegraphics[width=\textwidth]{C:/Users/Sina.Mokhtar.XLSCIENTIFIC/Documents/Problems/Copula_AD/figures/AC75a-5_DOB_072519_TS_2020-04-06_17_30_04_allCh.png}
		%     \caption{EEG plot from 2020-04-06 17:30:04 to 2020-04-07 17:30:04.}
		% \end{figure}
	
	\subsection{Zoomed-In Visualizations Before Seizures (March 25, 2020)}
	
	This section shows EEG segments before the first two seizures on March 25, 2020. Three time intervals were selected: 60 seconds, 1 hour (3600 seconds), and 100 minutes (6000 seconds) before seizure onset. Each plot includes the seizure event and the selected time range before it.
	
	\paragraph{First Seizure}
	
	\begin{figure}[h!]
		\centering
		\includegraphics[width=\textwidth, height=6cm]{C:/Users/Sina.Mokhtar.XLSCIENTIFIC/Documents/Problems/Copula_AD/figures/march25_60/AC75a-5_DOB 072519_TS_2020-03-25_17_30_04_seizure_00.png}
		\caption{EEG segment starting 60 seconds before the first seizure on March 25, 2020.}
	\end{figure}
	
	 \begin{figure}[h!]
		     \centering
		     \includegraphics[width=\textwidth, height=6cm]{C:/Users/Sina.Mokhtar.XLSCIENTIFIC/Documents/Problems/Copula_AD/figures/march25_600/AC75a-5_DOB 072519_TS_2020-03-25_17_30_04_seizure_00.png}
		     \caption{EEG segment starting 1 hour before the first seizure on March 25, 2020.}
		 \end{figure}
	
	\begin{figure}[h!]
		\centering
		\includegraphics[width=\textwidth, height=6cm]{C:/Users/Sina.Mokhtar.XLSCIENTIFIC/Documents/Problems/Copula_AD/figures/march25_3600/AC75a-5_DOB 072519_TS_2020-03-25_17_30_04_seizure_00.png}
		\caption{EEG segment starting 100 minutes before the first seizure on March 25, 2020.}
	\end{figure}
	
	
	
	\newpage
	
	
	
	\paragraph{Second Seizure}
	
	
	\begin{figure}[h!]
		\centering
		\includegraphics[width=\textwidth, height=6cm]{C:/Users/Sina.Mokhtar.XLSCIENTIFIC/Documents/Problems/Copula_AD/figures/march25_60/AC75a-5_DOB 072519_TS_2020-03-25_17_30_04_seizure_01.png}
		\caption{EEG segment starting 60 seconds before the second seizure on March 25, 2020.}
	\end{figure}
	
	 \begin{figure}[h!]
		     \centering
		     \includegraphics[width=\textwidth, height=6cm]{C:/Users/Sina.Mokhtar.XLSCIENTIFIC/Documents/Problems/Copula_AD/figures/march25_600/AC75a-5_DOB 072519_TS_2020-03-25_17_30_04_seizure_01.png}
		     \caption{EEG segment starting 1 hour before the second seizure on March 25, 2020.}
		 \end{figure}
	
	 \begin{figure}[h!]
		     \centering
		\includegraphics[width=\textwidth, height=6cm]{C:/Users/Sina.Mokhtar.XLSCIENTIFIC/Documents/Problems/Copula_AD/figures/march25_3600/AC75a-5_DOB 072519_TS_2020-03-25_17_30_04_seizure_01.png}
		\caption{EEG segment starting 100 minutes before the second seizure on March 25, 2020.}
	\end{figure}
	
	
	\section{Anomaly Detection Using Joint CDFs}
	
	Anomaly detection was performed using residual Variance Anomaly Detector (VAD) computed over overlapping windows of the EEG signal. A sliding window of 2000 samples (1 second at 2000 Hz) with an overlap of 500 samples was applied to all four channels. For each window, $\sigma^2$ was estimated using ARIMA model of order (1, 1, 1).
	
	Copula-based methods were used to model the joint CDF distribution of $\sigma^2$ across EEG channels. Both copulas transform marginal distributions into a joint cumulative distribution function (CDF), allowing anomaly scoring based on joint rarity.
	
	\begin{figure}[h!]
		\centering
		\includegraphics[width=\textwidth,height=8cm]{C:/Users/Sina.Mokhtar.XLSCIENTIFIC/Documents/Problems/Copula_AD/figures/sigma2_jointcdf_with_seizures_gumbel.png}
		\caption{Joint CDF series using Gaussian model.}
	\end{figure}
	
	\begin{figure}[h!]
		\centering
		\includegraphics[width=\textwidth,height=8cm]{C:/Users/Sina.Mokhtar.XLSCIENTIFIC/Documents/Problems/Copula_AD/figures/sigma2_jointcdf_with_seizures.png}
		\caption{Joint CDF series using Gumbel model.}
	\end{figure}
	
	
	
	\section{Anomaly Detection Using Autoencoder Reconstruction Loss}
	
	A autoencoder was trained using only non-seizure EEG windows. Each window was 600,000 samples long (5 minutes) and extracted with a step size of 20,000 samples. To ensure clean training data, windows within a 2-hour buffer around each seizure were excluded. During training, all data were downsampled by a factor of 100.
	
	The trained model was evaluated on separate EEG files containing seizures. For each test window, the reconstruction loss (mean squared error) was computed. Higher loss values indicate deviations from normal patterns and may reflect seizure activity or pre-ictal changes.
	
	The following figures show reconstruction loss over time, with seizure periods marked in red and the 1-hour pre-seizure windows shaded in orange.
	
	
\begin{figure}[h!]
	\centering
	\includegraphics[width=\textwidth]{../figures/reconstruction_loss_0330}
	\caption{Reconstruction loss with seizure intervals (March 30, 2020).}
	\label{fig:reconstructionloss0330}
\end{figure}
	
	
\end{document}
